\chapter{그래서 무엇을 알게 되었을까}

긴 계단을 다 올라왔어요.
교실에서 지우개를 빌리는 이야기로 시작해서, 블록체인 공책을 만들고,
점과 화살표로 지도를 그리고, 쪽지를 돌리고, 미래를 맞히는 시험까지 보았어요.
이제 위에서 내려다보며 정리해 볼게요.

\section{알게 된 것}

\begin{enumerate}
\item \textbf{허락 기록으로도 평판 점수를 만들 수 있어요.}
      “써도 돼”라는 허락 줄로 지도를 그리고 쪽지 돌리기 놀이를 하면 엔도스랭크 점수가 나와요.
\item \textbf{엔도스랭크와 송금 랭크는 서로 다른 것을 재요.}
      같은 지갑들을 두 점수로 줄 세우면 $\tau$가 0.360밖에 되지 않아요.
      허락 검사에서는 엔도스랭크가, 송금 검사에서는 송금 랭크가 더 잘 맞아요.
\item \textbf{엔도스랭크는 훨씬 빨라요.}
      화살표가 9배쯤 적어서, 같은 지갑들의 점수를 9배쯤 빨리 계산해요.
\item \textbf{미래를 맞힐 때도 각자 자기 쪽을 더 잘 맞혀요.}
      누가 새 허락을 받을지는 엔도스랭크가, 누가 새 송금자를 맞을지는 송금 랭크가 더 잘 맞혔어요.
      봄(3\textasciitilde{}5월)과 여름(6\textasciitilde{}8월) 두 번 모두 그랬어요.
      다만 엔도스랭크가 크게 앞선 까닭의 상당 부분은, 송금 랭크가 송금 기록이 없는 통에 대해 아무 말도 할 수 없다는 데 있었어요.
\item \textbf{하지만 받은 화살표를 세기만 해도 더 잘 맞혀요.}
      이 질문에서는 단순한 방법이 가장 강한 맞수였어요.
\item \textbf{두 지도를 섞은 동전 랭크는 미리 정한 약속을 지키지 못했어요.}
      9월 14일의 약속도, 10월 1일에 등록한 약속도요.
      새 허락을 맞히는 데에서는 송금 지도만 걷는 뒷면 랭크보다 확실히 나았지만, 새 송금자를 맞히는 데에서는 그렇지 못했어요.
      송금 랭크와 견주면 새 송금자에서 분명히 뒤졌어요.
\end{enumerate}

\section{주장하지 않는 것}

과학에서는 “무엇을 알았다”만큼 “무엇은 모른다”를 분명히 적는 것도 중요해요.
저는 논문에 이런 것들은 주장하지 \emph{않는다}고 적었어요.

\begin{itemize}
\item 송금 랭크가 이제 쓸모없다는 것.
\item 허락과 송금을 섞는 어떤 방법도 효과가 없다는 것. 제가 시험한 한 가지 섞는 방법이 약속을 지키지 못했을 뿐이에요.
\item 동전 랭크가 송금 지도만 걷는 방법이나 송금 랭크보다 낫다는 것.
\item 엔도스랭크가 새 허락을 잘 맞힌 것이 송금 기록이 없는 통과 상관없이도 그렇다는 것.
\item 엔도스랭크나 송금 랭크가 받은 화살표를 세는 것보다 미래를 잘 맞힌다는 것.
\item 제가 정한 숫자들(신선도가 줄어드는 빠르기, 양을 줄이는 정도)이나, 엔도스랭크에서 선생님이 쪽지를 모두에게 똑같이 주는 규칙이 가장 좋은 선택이라는 것.
\item 지도가 훨씬 커져도 엔도스랭크가 늘 그만큼 빠르다는 것.
\item 다른 블록체인이나 다른 기간에서도 같은 결과가 나온다는 것.
\item 이런 점수가 담보나 신분증을 대신할 수 있다는 것.
\end{itemize}

\section{아직 남은 질문}

이 연구가 끝난 자리에서 새 질문들이 시작돼요.

\begin{itemize}
\item 다른 블록체인에서도 허락 지도가 같은 이야기를 들려줄까요?
\item 지갑이 훨씬 많은 큰 지도에서도 엔도스랭크가 빠를까요? 확장 지갑 풀의 허락과 송금 기록을 끝까지 가져오는 일은 비용 때문에 그만두었는데, 다음 연구에서 마무리하고 싶어요.
\item 허락과 송금을 더 잘 섞는 방법이 있을까요? 동전 대신 지갑마다 다른 비율을 쓰면 어떨까요?
\item 보낸 양과 허락한 양을 “한 번”으로 줄이지 않고, 토큰의 값어치를 맞추어 쓰면 어떨까요? 봄에는 양을 그대로 쓴 앞면 랭크가 엔도스랭크보다 새 허락을 잘 맞혔거든요.
\item 단순한 세기가 이렇게 강하다면, 세기와 쪽지 놀이를 함께 쓰는 방법은 어떨까요?
\end{itemize}

\section{마지막으로}

이 책을 쓰면서 가장 전하고 싶었던 것은 숫자가 아니에요.
과학이 움직이는 방식이에요.

궁금한 것이 생기면 작은 것부터 차근차근 쌓아 올려요.
결과를 보기 전에 무엇이 성공인지 먼저 적어 둬요.
결과가 마음에 들지 않아도 그대로 알려요.
그리고 모르는 것은 모른다고 말해요.

여러분도 교실에서 “누구를 믿을까?”를 생각할 때가 있을 거예요.
그때 이 책의 쪽지 놀이를 떠올려 보세요.
누가 나를 믿는지뿐 아니라, \emph{나를 믿어 주는 친구를 누가 믿는지}도 중요하다는 것을요.
그리고 내가 열어 둔 문 하나하나가 누군가에게는 “나는 너를 믿어”라는 말이 된다는 것도요.

끝까지 읽어 주어서 고마워요.
